% =============================================================
% Bilgisayar Mimarisi Kitabı - Preamble
% Oğuz Ergin
% Şablon C: Modern Teknik (Libertinus + renkli başlıklar)
% =============================================================

% --- Temel Paketler ---
\usepackage{fontspec}
\usepackage[turkish,shorthands=:!]{babel}

% --- Yazı Tipi: Libertinus ---
\usepackage{libertinus}
% Monospace: Consolas (Windows) veya DejaVu Sans Mono (Linux/CI)
\IfFontExistsTF{Consolas}{%
  \setmonofont{Consolas}[Scale=0.82]%
}{%
  \setmonofont{DejaVu Sans Mono}[Scale=0.82]%
}

% --- Sayfa Düzeni ---
\usepackage[a4paper, margin=2.5cm, top=3cm, bottom=3cm, headheight=25pt]{geometry}
\raggedbottom  % Sayfaları aynı yüksekliğe zorlamayı kapat (aşırı dikey boşluk önlenir)

% --- Renkler (font/listing/titlesec tanımlarından önce) ---
\usepackage{xcolor}
\usepackage{pifont}  % \ding{51} (checkmark) ve \ding{55} (cross) için
\definecolor{kuturenk}{RGB}{230, 240, 250}
\definecolor{baslikrenk}{RGB}{0, 71, 133}
\definecolor{bolumrenk}{RGB}{0, 71, 133}
\definecolor{acikgri}{RGB}{245, 245, 245}
\definecolor{yanilgirenk}{RGB}{180, 40, 40}   % koyu kırmızı
\definecolor{tuzakrenk}{RGB}{180, 120, 0}      % koyu turuncu

% Yanılgı ve Tuzak başlık komutları
\newcommand{\yanilgi}[1]{\subsubsection*{{\color{yanilgirenk}\textbf{Yanılgı:}} #1}}
\newcommand{\tuzak}[1]{\subsubsection*{{\color{tuzakrenk}\textbf{Tuzak:}} #1}}

% TikZ blok şeması renkleri tikz-sekiller-makrolar.tex'te tanımlıdır.

% --- Bölüm Başlık Stili ---
\usepackage{titlesec}

% Bölüm başlığı: Renkli büyük numara + alt çizgi
\titleformat{\chapter}[display]
    {\normalfont\bfseries\color{bolumrenk}}
    {\flushright\fontsize{72}{72}\selectfont\color{bolumrenk!30}\thechapter}
    {-20pt}
    {\huge\flushleft}
    [\vspace{-5pt}{\color{bolumrenk}\titlerule[2pt]}]

% Section: renkli
\titleformat{\section}
    {\normalfont\Large\bfseries\color{bolumrenk}}
    {\thesection}{1em}{}

% Subsection
\titleformat{\subsection}
    {\normalfont\large\bfseries\color{bolumrenk!80}}
    {\thesubsection}{1em}{}
\titlespacing*{\subsection}{0pt}{1.2ex plus 0.4ex minus 0.2ex}{0.4ex plus 0.1ex}

% --- Matematik ---
\usepackage{amsmath}
\usepackage{unicode-math}
\setmathfont{Libertinus Math}

% Checkmark: Libertinus'ta ✓ glifi yok, platformda bulunan fonttan al
\IfFontExistsTF{Apple Symbols}{%
  \newfontfamily\symbolfont{Apple Symbols}%
}{%
  \IfFontExistsTF{Segoe UI Symbol}{%
    \newfontfamily\symbolfont{Segoe UI Symbol}%
  }{%
    \newfontfamily\symbolfont{DejaVu Sans}%
  }%
}
\AtBeginDocument{\def\checkmark{{\symbolfont ✓}}}

% Emoji desteği (Segoe UI Emoji veya Noto Color Emoji)
\IfFontExistsTF{Segoe UI Emoji}{%
  \newfontface\emojifont{Segoe UI Emoji}%
}{%
  \IfFontExistsTF{Noto Color Emoji}{%
    \newfontface\emojifont{Noto Color Emoji}%
  }{%
    \newfontface\emojifont{Segoe UI Symbol}%  % son çare: siyah-beyaz
  }%
}
\newcommand{\emoji}[1]{{\emojifont #1}}

% Kızgın Kuş ikonu (satır içi TikZ — orijinal çizim)
\newcommand{\kizginkus}{%
  \begin{tikzpicture}[baseline=-0.5ex, scale=0.12]
    % Gövde (kırmızı daire)
    \fill[red!85!black] (0,0) circle (1);
    % Karın (açık krem)
    \fill[orange!20!white] (0,-0.35) ellipse (0.55 and 0.4);
    % Sol göz (beyaz)
    \fill[white] (-0.3,0.2) ellipse (0.25 and 0.28);
    % Sağ göz (beyaz)
    \fill[white] (0.3,0.2) ellipse (0.25 and 0.28);
    % Sol göz bebeği (siyah, içe bakan)
    \fill[black] (-0.18,0.18) circle (0.1);
    % Sağ göz bebeği (siyah, içe bakan)
    \fill[black] (0.18,0.18) circle (0.1);
    % Sol kızgın kaş
    \fill[black] (-0.55,0.55) -- (-0.1,0.4) -- (-0.12,0.48) -- cycle;
    % Sağ kızgın kaş
    \fill[black] (0.55,0.55) -- (0.1,0.4) -- (0.12,0.48) -- cycle;
    % Gaga (sarı üçgen)
    \fill[yellow!90!orange] (0,-0.05) -- (-0.18,-0.2) -- (0.18,-0.2) -- cycle;
    % Tepedeki tüyler
    \fill[red!70!black] (-0.05,0.95) -- (0.05,0.95) -- (0.1,1.35) -- (0,1.25) -- (-0.1,1.35) -- cycle;
  \end{tikzpicture}%
}

% --- Tablolar ---
\usepackage{booktabs}
\usepackage{array}
\usepackage{multirow}
\usepackage{longtable}
\usepackage{tabularx}

% --- Listeler ---
\usepackage[shortlabels]{enumitem}
\setlist{itemsep=2pt, parsep=0pt, topsep=4pt}  % Liste maddelerindeki aşırı boşluğu azalt

% --- Şekiller ve Grafikler ---
\usepackage{graphicx}
\usepackage{float}
\usepackage{subcaption}
\captionsetup{labelfont=bf}                    % "Tablo 1.1:" ve "Şekil 1.1:" kalın
\captionsetup[table]{textfont=it}              % Tablo açıklama metni yatık
% Şekil/tablo etrafı boşluk (varsayılan ~20pt, sıkıştırılıyor)
\setlength{\textfloatsep}{10pt plus 3pt minus 2pt}
\setlength{\floatsep}{8pt plus 2pt minus 2pt}
\setlength{\intextsep}{8pt plus 2pt minus 2pt}
% Sayfa sonu boşluklarını azalt: float'lara daha fazla alan ver
\renewcommand{\topfraction}{0.9}        % Sayfanın en fazla %90'ı üstte float olabilir
\renewcommand{\bottomfraction}{0.9}     % Sayfanın en fazla %90'ı altta float olabilir
\renewcommand{\textfraction}{0.1}       % Sayfanın en az %10'u metin olmalı (varsayılan %20)
\renewcommand{\floatpagefraction}{0.75} % Sadece float sayfası için en az %75 dolu olmalı
\setcounter{topnumber}{3}               % Sayfa başına en fazla 3 üst float
\setcounter{bottomnumber}{2}            % Sayfa başına en fazla 2 alt float
\setcounter{totalnumber}{5}             % Sayfa başına toplam en fazla 5 float
\usepackage{tikz}
\usepackage{pgfplots}
\pgfplotsset{compat=1.18}
\usetikzlibrary{shapes, arrows, arrows.meta, positioning, calc, decorations.markings,
    circuits.logic.US, circuits.logic.IEC, fit, patterns, matrix, automata, babel,
    decorations.pathreplacing, decorations.pathmorphing, shadows, backgrounds}
\usepackage[section]{placeins}  % Şekillerin bölüm dışına taşmasını engeller
\usepackage[american]{circuitikz}  % Mantık kapısı devre şemaları

% TikZ global ok stili — tüm tikzpicture'larda >=Stealth varsayılan olur
\tikzset{>=Stealth}
\tikzset{okstil/.style={->, >=Stealth, thick}}

% --- TikZ Şekil Makroları ---
% =============================================================
% TikZ Şekil Makroları - Bilgisayar Mimarisi Kitabı
% Oğuz Ergin
% =============================================================
% Bu dosya tüm bölümlerde kullanılan ortak TikZ şekil
% makrolarını ve stil tanımlarını içerir.
% =============================================================

% -----------------------------------------------
% RENKLER
% -----------------------------------------------
\definecolor{sinyalmavi}{RGB}{70, 130, 180}
\definecolor{sinyalkirmizi}{RGB}{200, 60, 60}
\definecolor{sinyalyesil}{RGB}{60, 160, 60}
\definecolor{blokgri}{RGB}{240, 240, 240}
\definecolor{blokmavi}{RGB}{200, 220, 240}
\definecolor{blokkirmizi}{RGB}{255, 220, 220}
\definecolor{blokyesil}{RGB}{220, 245, 220}
\definecolor{bloksari}{RGB}{255, 250, 205}
\definecolor{blokturuncu}{RGB}{255, 225, 180}
\definecolor{vurgumavi}{RGB}{0, 100, 200}
\definecolor{vurgukirmizi}{RGB}{200, 0, 0}
\definecolor{vurguyesil}{RGB}{0, 140, 60}
\definecolor{vurguturuncu}{RGB}{200, 120, 0}

% -----------------------------------------------
% RENK VARYANTLARI (açık tonlar)
% -----------------------------------------------
\colorlet{blokmavi-acik}{blokmavi!30}
\colorlet{blokkirmizi-acik}{blokkirmizi!30}
\colorlet{blokyesil-acik}{blokyesil!30}
\colorlet{bloksari-acik}{bloksari!30}
\colorlet{blokgri-acik}{blokgri!30}

% -----------------------------------------------
% BORU HATTI AŞAMA RENKLERİ
% -----------------------------------------------
\definecolor{asamagetir}{RGB}{200, 220, 240}    % G: Getir — mavi
\definecolor{asamacoz}{RGB}{220, 245, 220}       % Ç: Çöz — yeşil
\definecolor{asamayurut}{RGB}{255, 235, 200}     % U: Yürüt — turuncu
\definecolor{asamabellek}{RGB}{255, 220, 220}    % B: Bellek — kırmızı
\definecolor{asamayaz}{RGB}{230, 220, 250}       % Y: Yaz — mor

% -----------------------------------------------
% GENEL TikZ STİLLERİ
% -----------------------------------------------

% Blok şeması stilleri
\tikzset{
    % Temel blok stili
    blok/.style={
        draw, rectangle, rounded corners=2pt,
        minimum width=2cm, minimum height=1cm,
        fill=blokmavi, font=\small,
        align=center
    },
    % Vurgulu blok
    vurgublok/.style={
        blok, fill=blokkirmizi, draw=vurgukirmizi, thick
    },
    % Bellek bloku
    bellekblok/.style={
        draw, rectangle, minimum width=2.5cm, minimum height=1.5cm,
        fill=blokyesil, font=\small, align=center
    },
    % Yazmaç bloku
    yazmacblok/.style={
        draw, rectangle, minimum width=2cm, minimum height=0.8cm,
        fill=bloksari, font=\small, align=center
    },
    % Ok stilleri
    okstil/.style={
        ->, >=Stealth, thick
    },
    ciftok/.style={
        <->, >=Stealth, thick
    },
    veristil/.style={
        ->, >=Stealth, thick, sinyalmavi
    },
    denetimstil/.style={
        ->, >=Stealth, dashed, sinyalkirmizi
    },
}

% -----------------------------------------------
% KARNO HARİTASI MAKROSU
% -----------------------------------------------

% 2 değişkenli Karno haritası
% #1: değişken isimleri (örn. "A" ve "B")
% #2: hücre değerleri (4 adet, sırasıyla 00,01,10,11)
\newcommand{\karnoiki}[4][]{%
    % #2=değişken1, #3=değişken2, #4=değerler (virgülle ayrılmış 4 değer)
    \begin{tikzpicture}[scale=0.8]
        % Grid
        \draw (0,0) grid (2,2);
        % Değişken etiketleri
        \node[above] at (0.5, 2) {\footnotesize 0};
        \node[above] at (1.5, 2) {\footnotesize 1};
        \node[left] at (0, 1.5) {\footnotesize 0};
        \node[left] at (0, 0.5) {\footnotesize 1};
        \node[above left] at (0, 2) {\footnotesize #2\textbackslash #3};
        % Hücre değerleri
        \foreach \val [count=\i from 0] in {#4} {
            \pgfmathtruncatemacro{\col}{mod(\i,2)}
            \pgfmathtruncatemacro{\row}{1 - int(\i/2)}
            \node at (\col+0.5, \row+0.5) {\val};
        }
        #1
    \end{tikzpicture}
}

% 3 değişkenli Karno haritası
% #1: ek TikZ kodu (gruplamalar vb.)
% #2: satır değişkeni, #3: sütun değişkenleri
% #4: 8 hücre değeri (AB\C sıralaması: 00/0, 00/1, 01/0, 01/1, 11/0, 11/1, 10/0, 10/1)
\newcommand{\karnouc}[4][]{%
    \begin{tikzpicture}[scale=0.8]
        \draw (0,0) grid (4,2);
        \node[above] at (0.5, 2) {\footnotesize 00};
        \node[above] at (1.5, 2) {\footnotesize 01};
        \node[above] at (2.5, 2) {\footnotesize 11};
        \node[above] at (3.5, 2) {\footnotesize 10};
        \node[left] at (0, 1.5) {\footnotesize 0};
        \node[left] at (0, 0.5) {\footnotesize 1};
        \node[above left] at (0, 2) {\footnotesize #2\textbackslash #3};
        \foreach \val [count=\i from 0] in {#4} {
            \pgfmathtruncatemacro{\col}{mod(\i,4)}
            \pgfmathtruncatemacro{\row}{1 - int(\i/4)}
            \node at (\col+0.5, \row+0.5) {\val};
        }
        #1
    \end{tikzpicture}
}

% 4 değişkenli Karno haritası
% #1: ek TikZ kodu
% #2: satır değişkenleri, #3: sütun değişkenleri
% #4: 16 hücre değeri (AB\CD sıralaması: Karno sırasıyla 00,01,11,10)
\newcommand{\karnodort}[4][]{%
    \begin{tikzpicture}[scale=0.8]
        \draw (0,0) grid (4,4);
        \node[above] at (0.5, 4) {\footnotesize 00};
        \node[above] at (1.5, 4) {\footnotesize 01};
        \node[above] at (2.5, 4) {\footnotesize 11};
        \node[above] at (3.5, 4) {\footnotesize 10};
        \node[left] at (0, 3.5) {\footnotesize 00};
        \node[left] at (0, 2.5) {\footnotesize 01};
        \node[left] at (0, 1.5) {\footnotesize 11};
        \node[left] at (0, 0.5) {\footnotesize 10};
        \node[above left] at (0, 4) {\footnotesize #2\textbackslash #3};
        \foreach \val [count=\i from 0] in {#4} {
            \pgfmathtruncatemacro{\col}{mod(\i,4)}
            \pgfmathtruncatemacro{\row}{3 - int(\i/4)}
            \node at (\col+0.5, \row+0.5) {\val};
        }
        #1
    \end{tikzpicture}
}

% -----------------------------------------------
% FLIP-FLOP MAKROLARI
% -----------------------------------------------

% Genel flip-flop çizim stili
\tikzset{
    flipflop/.style={
        draw, rectangle, minimum width=1.8cm, minimum height=2.4cm,
        font=\small, thick
    },
    ffgiris/.style={font=\footnotesize},
    ffcikis/.style={font=\footnotesize},
}

% -----------------------------------------------
% MANTIK KAPISI STİLLERİ
% -----------------------------------------------

\tikzset{
    % Mantık kapısı sembol boyutları
    kapiboy/.style={
        circuit logic US, huge circuit symbols
    },
}

% -----------------------------------------------
% BIT ALANI DİYAGRAMI MAKROSU
% -----------------------------------------------

% Buyruk bit alanı çizimi
% #1: toplam bit genişliği (örn. 32)
% #2: alanlar listesi (tikz node'ları olarak verilecek)
\tikzset{
    bitalani/.style={
        draw, minimum height=0.8cm, font=\footnotesize\ttfamily,
        anchor=west, fill=blokmavi
    },
    bitust/.style={
        font=\tiny, above, text=gray
    },
    bitalt/.style={
        font=\scriptsize, below
    },
}

% -----------------------------------------------
% AKIŞ ŞEMASI STİLLERİ
% -----------------------------------------------

\tikzset{
    % Akış şeması blokları
    akisbaslangic/.style={
        draw, rounded rectangle, rounded rectangle west arc=10pt,
        rounded rectangle east arc=10pt,
        minimum width=2cm, minimum height=0.6cm,
        fill=blokyesil, font=\footnotesize, align=center,
        thick
    },
    akisislem/.style={
        draw, rectangle, rounded corners=2pt,
        minimum width=2.5cm, minimum height=0.6cm,
        fill=blokmavi, font=\footnotesize, align=center,
        thick
    },
    akiskarar/.style={
        draw, diamond, aspect=2, minimum width=1.8cm,
        fill=bloksari, font=\footnotesize, align=center,
        inner sep=1pt, thick
    },
    akisok/.style={
        ->, >=Stealth, thick
    },
    % Geri dönüş oku (döngü çizgileri için)
    akisdonus/.style={
        ->, >=Stealth, thick, rounded corners=8pt
    },
}

% -----------------------------------------------
% ZAMANLAMA ŞEMASI STİLLERİ
% -----------------------------------------------

\tikzset{
    zamansinyal/.style={
        very thick, sinyalmavi
    },
    zamanok/.style={
        <->, >=Stealth, thin
    },
    zamanetiket/.style={
        font=\footnotesize
    },
}

% -----------------------------------------------
% VERİYOLU BİLEŞEN STİLLERİ (Bölüm 4-5 için)
% -----------------------------------------------

\tikzset{
    % Çoklayıcı (MUX)
    mux/.style={
        draw, trapezium, trapezium angle=70,
        minimum width=0.6cm, minimum height=1.2cm,
        font=\tiny, thick, fill=blokgri, shape border rotate=270
    },
    % AMB
    amb/.style={
        draw, rectangle, minimum width=1.2cm, minimum height=1.8cm,
        fill=blokmavi, font=\small, thick, align=center
    },
    % Toplayıcı (daire)
    toplayici/.style={
        draw, circle, minimum size=0.6cm, font=\small, thick,
        fill=blokgri, inner sep=0pt
    },
    % Yazmaç (ince uzun dikdörtgen)
    yazmac/.style={
        draw, rectangle, minimum width=0.3cm, minimum height=2cm,
        fill=bloksari, font=\tiny, thick
    },
    % Denetim birimi
    denetim/.style={
        draw, ellipse, minimum width=2cm, minimum height=1.2cm,
        fill=blokkirmizi!30, font=\small, thick, align=center
    },
    % Bellek
    bellek/.style={
        draw, rectangle, minimum width=1.8cm, minimum height=2cm,
        fill=blokyesil, font=\small, thick, align=center
    },
    % Veri yolu (kalın çizgi)
    veriyolu/.style={
        line width=2pt, ->, >=Stealth
    },
    % İnce sinyal
    sinyal/.style={
        ->, >=Stealth, thin
    },
}

% -----------------------------------------------
% DURUM MAKİNESİ STİLLERİ
% -----------------------------------------------

\tikzset{
    durum/.style={
        draw, circle, minimum size=1.2cm, font=\small,
        thick, fill=blokmavi
    },
    durumok/.style={
        ->, >=Stealth, thick, bend left=20
    },
    durumdongu/.style={
        ->, >=Stealth, thick, loop above, looseness=8
    },
}

% -----------------------------------------------
% BELLEK DİYAGRAMI STİLLERİ (Bölüm 6 için)
% -----------------------------------------------

\tikzset{
    % Önbellek satırı
    onbellek/.style={
        draw, rectangle, minimum height=0.5cm, font=\tiny,
        fill=blokmavi!30
    },
    % Bellek hücresi
    bellekhucre/.style={
        draw, rectangle, minimum width=0.8cm, minimum height=0.5cm,
        font=\tiny
    },
}

% -----------------------------------------------
% BORU HATTI DİYAGRAM STİLLERİ (Bölüm 5-6 için)
% -----------------------------------------------

\tikzset{
    % Pipeline aşama kutusu
    asama/.style={
        draw, rectangle, minimum width=1.2cm, minimum height=0.7cm,
        font=\footnotesize, align=center, thick
    },
    % Beş aşama stili
    asamaG/.style={asama, fill=asamagetir},
    asamaC/.style={asama, fill=asamacoz},
    asamaU/.style={asama, fill=asamayurut},
    asamaB/.style={asama, fill=asamabellek},
    asamaY/.style={asama, fill=asamayaz},
    % Boş/stall aşama
    asamabos/.style={asama, fill=blokgri, dashed},
    % Pipeline yazmaç (aşamalar arası)
    pipeyazmac/.style={
        draw, rectangle, minimum width=0.15cm, minimum height=1.8cm,
        fill=bloksari, thick
    },
    % İşlemci birimi (veriyolu bileşeni)
    birim/.style={
        draw, rectangle, rounded corners=3pt,
        minimum width=1.5cm, minimum height=1cm,
        fill=blokmavi, font=\small, thick, align=center
    },
}

% -----------------------------------------------
% ÇOK ÇEKİRDEKLİ / PARALEL STİLLER (Bölüm 8-9)
% -----------------------------------------------

\tikzset{
    % İşlemci çekirdeği
    cekirdek/.style={
        draw, rectangle, rounded corners=4pt,
        minimum width=2cm, minimum height=1.5cm,
        fill=blokmavi, font=\small, thick, align=center
    },
    % Önbellek kutusu
    onbellekblok/.style={
        draw, rectangle, minimum width=1.8cm, minimum height=0.8cm,
        fill=blokyesil-acik, font=\footnotesize, align=center
    },
    % Paylaşılan kaynak
    paylasilan/.style={
        draw, rectangle, rounded corners=2pt,
        minimum width=3cm, minimum height=1cm,
        fill=bloksari, font=\small, thick, align=center
    },
    % Bağlantı yolu (bus/interconnect)
    baglanti/.style={
        draw, rectangle, minimum height=0.4cm,
        fill=blokgri, font=\footnotesize
    },
}

% -----------------------------------------------
% GPU / HIZLANDIRICI STİLLERİ (Bölüm 10)
% -----------------------------------------------

\tikzset{
    % SM / CU birimi
    gpubirim/.style={
        draw, rectangle, rounded corners=2pt,
        minimum width=1cm, minimum height=0.8cm,
        fill=blokmavi, font=\tiny, thick, align=center
    },
    % Bellek modülü
    gpubellek/.style={
        draw, rectangle, minimum width=2cm, minimum height=0.6cm,
        fill=blokyesil, font=\footnotesize, align=center
    },
}


% --- Kod Listeleme ---
\usepackage{listings}
\renewcommand{\lstlistingname}{Kod}
\lstloadlanguages{[ANSI]C}  % listings 1.11b uyumsuzluk düzeltmesi: C dilini önceden yükle
\lstset{
    basicstyle=\ttfamily\footnotesize,
    breaklines=true,
    frame=single,
    numbers=left,
    numberstyle=\tiny\color{gray},
    tabsize=4,
    captionpos=b,
    xleftmargin=2em,
    framexleftmargin=1.5em,
    backgroundcolor=\color{acikgri},
    rulecolor=\color{bolumrenk!40},
    extendedchars=true,
    columns=fullflexible,
    keepspaces=true,
    texcl=true,
}

% --- Assembly/MIPS stil tanımı ---
\lstdefinestyle{mips}{
    morekeywords={add, sub, addi, lw, sw, beq, bne, slt, j, jal, jr, and, or, sll, srl, lui, ori, mult, div, mfhi, mflo},
    keywordstyle=\color{blue}\bfseries,
    commentstyle=\color{green!50!black}\ttfamily,
    stringstyle=\color{red},
    morecomment=[l]{\#},
    texcl=true,
}

% --- Assembly/RISC-V stil tanımı ---
\lstdefinestyle{riscv}{
    morekeywords={add, sub, addi, and, or, xor, andi, ori, xori, sll, srl, sra, slli, srli, srai, slt, sltu, slti, sltiu, lw, sw, lb, sb, lh, sh, lbu, lhu, ld, sd, beq, bne, blt, bge, bltu, bgeu, jal, jalr, lui, auipc, ecall, ebreak, mul, mulh, mulhu, mulhsu, div, divu, rem, remu, nop, li, la, mv, not, neg, j, ret, call, tail, fadd.s, fsub.s, fmul.s, fdiv.s, fsqrt.s, fadd.d, fsub.d, fmul.d, fdiv.d, flw, fsw, fld, fsd, feq.s, flt.s, fle.s, fcvt.w.s, fcvt.s.w, csrr, csrw, csrs, csrc, fence, fence.i, wfi, mret, sret, lr.w, sc.w, lr.d, sc.d, amoswap.w, amoadd.w, amoand.w, amoor.w, amoxor.w, amomax.w, amomin.w, amomaxu.w, amominu.w, amoswap.d, amoadd.d, amoand.d, amoor.d, amoxor.d, amomax.d, amomin.d, amomaxu.d, amominu.d, prefetch.r, prefetch.w, prefetch.i, zero, ra, sp, gp, tp, t0, t1, t2, t3, t4, t5, t6, s0, s1, s2, s3, s4, s5, s6, s7, s8, s9, s10, s11, a0, a1, a2, a3, a4, a5, a6, a7, fp, x0, x1, x2, x3, x4, x5, x6, x7, x8, x9, x10, x11, x12, x13, x14, x15, x16, x17, x18, x19, x20, x21, x22, x23, x24, x25, x26, x27, x28, x29, x30, x31, f0, f1, f2, f3, f4, f5, f6, f7},
    keywordstyle=\color{blue}\bfseries,
    commentstyle=\color{green!50!black}\ttfamily,
    stringstyle=\color{red},
    morecomment=[l]{\#},
    texcl=true,
    sensitive=true,
}

% --- Verilog stil tanımı ---
\lstdefinestyle{verilog}{
    morekeywords={module, endmodule, input, output, inout, wire, reg, parameter, localparam, assign, always, begin, end, if, else, case, endcase, default, posedge, negedge, initial, forever, for, while, integer, genvar, generate, endgenerate, task, endtask, function, endfunction},
    keywordstyle=\color{blue}\bfseries,
    commentstyle=\color{green!50!black},
    stringstyle=\color{red},
    sensitive=true,
}

% --- CUDA/C stil tanımı ---
\lstdefinestyle{cuda}{
    morekeywords={__global__, __device__, __shared__, __host__, __constant__, int, void, if, else, return, for, while, float, double, const, unsigned, struct, typedef, sizeof, static, extern, char, short, long, signed, blockIdx, blockDim, threadIdx, gridDim, warpSize},
    keywordstyle=\color{blue}\bfseries,
    commentstyle=\color{green!50!black},
    stringstyle=\color{red},
}

% --- Alıştırmalardaki kod blokları için kompakt stil ---
\lstdefinestyle{alistirma-kod}{
    numbers=none,
    xleftmargin=1em,
    framexleftmargin=0.5em,
    frame=single,
    backgroundcolor=\color{acikgri},
    rulecolor=\color{bolumrenk!40},
    basicstyle=\ttfamily\small,
    breaklines=true,
    tabsize=4,
    aboveskip=0.5\baselineskip,
    belowskip=0.5\baselineskip,
}

% --- Bağlantılar ---
\usepackage{hyperref}
\hypersetup{
    colorlinks=true,
    linkcolor=bolumrenk,
    citecolor=bolumrenk,
    urlcolor=bolumrenk!70,
    bookmarks=true,
    bookmarksnumbered=true,
}

% --- Kaynakça ---
\usepackage{csquotes}
\DeclareQuoteAlias{american}{turkish}
\usepackage[style=numeric, sorting=nyt, backend=biber, defernumbers=false]{biblatex}
\addbibresource{kaynakca.bib}

% --- Dizin ---
\usepackage{makeidx}
\makeindex

% --- Özel Ortamlar (Environments) ---
\usepackage{tcolorbox}
\tcbuselibrary{breakable, skins}

% Örnek kutusu
\newtcolorbox{ornek}[1][]{
    colback=bolumrenk!5,
    colframe=bolumrenk,
    fonttitle=\bfseries,
    title={Örnek #1},
    breakable,
    enhanced,
    left=4mm,
    right=4mm,
    arc=1mm,
}

% Tanım kutusu
\newtcolorbox{tanim}[1][]{
    colback=orange!5!white,
    colframe=orange!80!black,
    fonttitle=\bfseries,
    title={#1},
    breakable,
    enhanced,
    arc=1mm,
}

% Alıştırma zorluk dereceleri
\newcommand{\kolay}{{\normalfont\small\textcolor{green!60!black}{★}}}
\newcommand{\orta}{{\normalfont\small\textcolor{orange!80!black}{★★}}}
\newcommand{\zor}{{\normalfont\small\textcolor{red!70!black}{★★★}}}

% Alıştırma ortamı — isteğe bağlı zorluk parametresi: \begin{alistirma}[\kolay]
\newcounter{alistirma}[chapter]
\renewcommand{\thealistirma}{\thechapter.\arabic{alistirma}}
\newenvironment{alistirma}[1][]{%
    \refstepcounter{alistirma}%
    \par\medskip\noindent\textbf{Alıştırma \thealistirma.}%
    \def\alistirmazorluk{#1}%
    \ifx\alistirmazorluk\empty\else\hspace{0.5em}#1\fi\space
}{%
    \par\medskip
}

% Bilgi notu sayacı ve listesi
\newcounter{bilginotu}[chapter]
\renewcommand{\thebilginotu}{\thechapter.\arabic{bilginotu}}

\makeatletter
\newcommand{\listofbilginotlari}{%
    \chapter*{Bilgi Notları Listesi}%
    \addcontentsline{toc}{chapter}{Bilgi Notları Listesi}%
    \@starttoc{lon}%
}
\newcommand{\l@bilginotu}{\@dottedtocline{1}{0em}{3em}}
\makeatother

% Bilgi notu iç kutusu (doğrudan kullanılmaz)
\newtcolorbox{bilginotuic}[1]{
    colback=green!5!white,
    colframe=green!60!black,
    fonttitle=\bfseries,
    title={#1},
    enhanced,
    breakable,
    arc=1mm,
    borderline west={3pt}{0pt}{green!60!black},
}

% Bilgi notu ortamı (numaralı, listeye eklenir)
\newenvironment{bilginot}[1][]{%
    \refstepcounter{bilginotu}%
    \def\bilginotarg{#1}%
    \ifx\bilginotarg\empty
        \addcontentsline{lon}{bilginotu}{\protect\numberline{\thebilginotu}Bilgi Notu}%
        \begin{bilginotuic}{Bilgi Notu~\thebilginotu}%
    \else
        \addcontentsline{lon}{bilginotu}{\protect\numberline{\thebilginotu}#1}%
        \begin{bilginotuic}{Bilgi Notu~\thebilginotu: #1}%
    \fi
}{%
    \end{bilginotuic}%
}

% Terim kökeni kutusu
\definecolor{terimrenk}{RGB}{139, 90, 43}
\newtcolorbox{terimnotu}[1][Terim Kökenleri]{
    colback=terimrenk!5!white,
    colframe=terimrenk!80!black,
    fonttitle=\bfseries,
    title={\small\textsf{#1}},
    breakable,
    enhanced,
    arc=1mm,
    borderline west={3pt}{0pt}{terimrenk!80!black},
    fontupper=\small,
    top=2mm, bottom=2mm,
}

% Karşılaştırma kutusu (diğer mimarilerle karşılaştırma)
\definecolor{karsirenk}{RGB}{100, 60, 150}
\newtcolorbox{karsilastirma}[1][Diğer Mimarilerde]{
    colback=karsirenk!5!white,
    colframe=karsirenk!70!black,
    fonttitle=\bfseries,
    title={\small\textsf{#1}},
    breakable,
    enhanced,
    arc=1mm,
    borderline west={3pt}{0pt}{karsirenk!70!black},
    fontupper=\small,
    top=2mm, bottom=2mm,
}

% --- Üstbilgi/Altbilgi ---
\usepackage{fancyhdr}
\pagestyle{fancy}
\fancyhf{}
\fancyhead[LE]{\small\color{bolumrenk}\leftmark}
\fancyhead[RO]{\small\color{bolumrenk}\rightmark}
\fancyfoot[C]{\color{bolumrenk}\thepage}
\renewcommand{\headrulewidth}{0.4pt}
\renewcommand{\headrule}{\hbox to\headwidth{\color{bolumrenk}\leaders\hrule height \headrulewidth\hfill}}
\renewcommand{\footrulewidth}{0pt}

% --- Şekil numaralama ---
\numberwithin{figure}{chapter}
\numberwithin{table}{chapter}
\numberwithin{equation}{chapter}

% --- Şekil/Tablo listesinde bölüm başlıkları ---
\makeatletter
\let\kitap@origchapter\@chapter
\def\@chapter[#1]#2{%
  \kitap@origchapter[#1]{#2}%
  \addtocontents{lof}{%
    {\protect\noindent\protect\bfseries\protect\color{bolumrenk}%
     \@chapapp~\thechapter{} --- #1\protect\par\protect\nopagebreak\protect\smallskip}%
  }%
  \addtocontents{lot}{%
    {\protect\noindent\protect\bfseries\protect\color{bolumrenk}%
     \@chapapp~\thechapter{} --- #1\protect\par\protect\nopagebreak\protect\smallskip}%
  }%
  \addtocontents{lon}{\protect\addvspace{10\p@}}%
  \addtocontents{lon}{%
    {\protect\noindent\protect\bfseries\protect\color{bolumrenk}%
     \@chapapp~\thechapter{} --- #1\protect\par\protect\nopagebreak\protect\smallskip}%
  }%
}
\makeatother

% --- Bölüm açılış fotoğrafı (başlıkla aynı sayfada) ---
\newcommand{\bolumfoto}[3][]{%
    \par\vspace{0.6em}%
    \begin{center}%
    \begin{minipage}{\textwidth}%
    \centering
    \includegraphics[width=0.82\textwidth,height=0.38\textheight,keepaspectratio]{#2}\\[6pt]%
    {\small\color{bolumrenk}\textit{#3}}%
    \ifx\relax#1\relax\else\\[8pt]%
    {\small\itshape #1}%
    \fi
    \end{minipage}%
    \end{center}%
    \par\vspace{0.8em}%
    \newpage%
}

% --- Eksik şekil yer tutucusu ---
\newcommand{\eksikfigur}[2][]{%
    \begin{figure}[H]
    \centering
    \fbox{\parbox{0.7\textwidth}{\centering\vspace{2cm}\textbf{[ŞEKİL GEREKLİ]}\\#2\vspace{2cm}}}
    \caption{#2}
    \ifx&#1&\else\label{#1}\fi
    \end{figure}
}
